\section{Introduction}
\label{sec:intro}

When \citet{bennett2009neural} placed a dead Atlantic salmon in an fMRI scanner and asked it to determine the emotions of photographs, the resulting ``significant'' brain activations became one of the most memorable demonstrations of the multiple comparisons problem in science.
The study showed that standard analytical pipelines could produce plausible but completely spurious results---a finding that reshaped neuroimaging methodology.
Today, a new class of analytical pipelines demands similar scrutiny: autonomous \aiscientist systems that generate hypotheses, analyze data, and write research papers with minimal human oversight \citep{lu2024ai}.

These systems are scaling rapidly.
\citet{lu2024ai} demonstrated end-to-end automated research at roughly \$15 per paper; subsequent work has identified numerous failure modes including coding errors \citep{beel2025evaluating}, fabricated results \citep{luo2025more}, and poor implementation capability \citep{zhu2025ai}.
Yet a fundamental question remains unanswered: \textbf{when given data with a known generating process, do \aiscientists correctly recover the true model?}
Existing evaluations either audit AI-generated papers post-hoc without ground truth \citep{luo2025more, beel2025evaluating} or benchmark replication ability on open-ended tasks \citep{starace2025paperbench}.
None provide a controlled experiment where the correct answer is known in advance.

We fill this gap with the \livesalmon---a complement to Bennett et al.'s original experiment.
Where the dead salmon test asks whether analytical pipelines produce false positives on null data, the live salmon test asks whether they recover true signal from structured data.
We generate a synthetic dataset from a known formula containing both linear terms ($X_1$, $X_2$, $X_3$) and a multiplicative interaction ($X_4 \times X_5$), plus correlated proxy features and pure noise dimensions.
Three frontier \llms (\gptmodel, \claudemodel, \geminimodel) each analyze this dataset three times, producing 9 ``live salmon'' papers.
We repeat the experiment with a null dataset (pure noise target), producing 9 ``dead salmon'' papers as controls.

Our results reveal a striking pattern.
All \aiscientists reliably detect the strong linear effects: $X_1$ (coefficient 2.0) and $X_3$ (coefficient $-1.5$) are identified in 100\% of papers, and the weaker $X_2$ (coefficient 0.5) in 67\%.
However, \textbf{zero out of 9 papers detect the $X_4 \times X_5$ interaction}, despite it contributing 8\% of total variance---comparable in effect size to $X_2$, which most papers do find.
On the dead salmon data, all 9 papers correctly report null findings with zero false positives, in direct contrast to Bennett et al.'s original result.
The collective ``AI scientific literature'' produced by these systems would lead readers to believe the data generating process is purely linear, systematically missing a meaningful nonlinear component.

We make the following contributions:
\begin{itemize}[leftmargin=*,itemsep=0pt,topsep=0pt]
    \item We introduce the \livesalmon, the first controlled benchmark that evaluates \aiscientists against a known data generating process, complementing the classic dead salmon paradigm.
    \item We conduct a systematic evaluation of three frontier \llms acting as data scientists, revealing that they universally fail to detect interaction effects (0/9) despite reliably identifying linear predictors (100\% for strong effects).
    \item We demonstrate that current \aiscientists exhibit zero false positives on null data---a surprising contrast to the original dead salmon fMRI study---revealing that the core limitation is not statistical conservatism but methodological narrowness.
    \item We provide a reusable evaluation framework with known-\dgp datasets and automated claim extraction for benchmarking future \aiscientist systems.
\end{itemize}
